\scanpage{210}
SYSTÈMES de RACINES (2) \hfill Tribu n$^\circ$ 59

Pour orienter le lecteur, nous indiquons ici les principaux exemples de systèmes de racines que l'on rencontre dans la théorie des groupes et algèbres de Lie.

I. Soient $\underline{g}$ une algèbre de Lie semi-simple sur $\underline{C}$ et $\underline{h}$ une sous-algèbre de Cartan de $\underline{g}$.
L'algèbre $\underline{h}$ est donc commutative, formée d'éléments semi-simples.
Une racine de $\underline{g}$ par rapport à $\underline{h}$ est une forme linéaire $\alpha$ non nulle sur $\underline{h}$ pour laquelle il existe $x\in\underline{g}$, $x\ne0$, tel que $[h,x]=\alpha(x)$ quel que soit $h\in\underline{h}$.
Soit $V$ l'espace vectoriel sur $\underline{Q}$ engendré par l'ensemble $R$ des racines de $\underline{g}$ par rapport à $\underline{h}$.
Alors $\underline{h}^*=V\otimes\underline{C}$, et $R$ est un système de racines réduit dans $V$, au sens du \S\ 6.
Tout système de racines réduit est l'ensemble des racines d'une algèbre semi-simple complexe, unique à son isomorphisme près.
Tout automorphisme de $\underline{g}$ laissant stable $\underline{h}$ induit un automorphisme de $\underline{h}^*$ laissant $R$ et $V$ invariants, et le groupe des automorphismes de $R$ ainsi obtenus à partir des automorphismes, (resp. des automorphismes intérieurs), de $\underline{g}$ s'identifie au groupe de tous les automorphismes de $R$, (resp. au groupe de Weyl de $R$).
Identifions le dual de $V$ à une $\underline{Q}$-forme $\underline{h}^0$ de $\underline{h}$ et soit $G$ un groupe de Lie complexe connexe d'algèbres de Lie $\underline{g}$.
Alors le noyau $\unsure{\Gamma|T}$ de l'application exponentielle de $\underline{h}$ dans $G$ est un réseau de $i.\underline{h}^0$ compris entre $2\pi i.Q(R^*)$ et $2\pi i.P(R^*)$, où $R^*$ est le système de racines inverse de $R$, qui est égal à $2\pi i.Q(R^*)$, (resp. $2\pi iP(R^*)$), si $G$ est simplement connexe, (resp. isomorphe à $\operatorname{Ad}\underline{g}$).
Le groupe $2\pi i.P(R^*)/\unsure{\Gamma|T}$ s'identifie au centre de $G$, et $\unsure{\Gamma|T}/2\pi i.Q(R^*)$ au groupe fondamental de $G$.
Le réseau de $V$ associé à $(2\pi i)^{-1}.\unsure{\Gamma|T}$ est engendré par les poids des représentations de $G$.

II. Soient $\underline{g}$ une algèbre de Lie semi-simple sur $\underline{R}$, $\underline{k}$ une sous-algèbre compacte maximale de $\underline{g}$, $\underline{p}$ le complément orthogonal de $\underline{k}$ par rapport à la forme de Killing, et $\underline{a}$ une sous-algèbre commutative maximale de $\underline{p}$. Pour toute forme

\margin{209}
